\documentclass[11pt,a4paper]{article}
\usepackage[T1]{fontenc}
\usepackage{lmodern,amsmath,amssymb,microtype,booktabs}
\usepackage[margin=25mm]{geometry}
\usepackage[colorlinks=true,urlcolor=blue,linkcolor=blue]{hyperref}
\newcommand{\Om}{\Omega}
\newcommand{\CC}{\mathcal C}
\title{Wilson-duality stabilization of the dS$_3$ crossing kernel}
\author{Alan Day \and Chengheng\thanks{AI research assistant (OpenAI).}}
\date{Working research note --- 31 October 2026}
\begin{document}\maketitle
\begin{abstract}
Groenevelt's Wilson-function duality converts the exponentially small
Omega factor in one shadow term into the complex conjugate of its
oppositely reflected partner, multiplied by an explicit Gamma ratio.
After combining coefficients, the two terms of the complete physical
scalar kernel have equal coefficient moduli. This removes the exponential
error amplification identified in our preceding note without discarding
phases. For one $M=32$ sample, a 128-term direct sum changes from an
absolute error near $5\,10^{43}$ to a discrepancy near $1.7\,10^{-42}$
against the saved reference. The new exact-arithmetic tail bound is
$1.1\,10^{-40}$. This is a stable representation and a pointwise
error bound, not a proof of a heavy-mass integral limit.
\end{abstract}

\section{The standard Wilson identity and our parameter dictionary}
The physical setting remains radius one, $G_N=0$, proper $SO^+(3,1)$,
BD scalar observers, KG matter, a fixed positive observer-label lower
bound, and half-period regulated local creation blocks.
External labels are $(a,q,c,p)$, internal labels $(b,d)>0$.
Write $h_i=(1+i\nu_i)/2$, $h=(1+ib)/2$, $h'=(1+id)/2$.
The scalar convention is $\bar h_i=h_i$, not conjugation.
We do not identify the scalar subblock with a complete unitary transform.

We use Groenevelt's meromorphic Wilson identities \cite{groenevelt},
rather than the positivity assumptions of his Hilbert-space transforms.
For $T=(T_1,\ldots,T_4)$ and $S=\sum_jT_j$, define the symmetric
$c$-function expansion
\begin{align}
 F(T;\lambda,x)
 &=\sum_{\epsilon=\pm1}
 \frac{\Gamma(-2i\epsilon\lambda)}
 {\prod_j\Gamma(T_j-i\epsilon\lambda)
  \Gamma(S/2+i\epsilon\lambda+ix)
  \Gamma(S/2+i\epsilon\lambda-ix)}
 \nonumber\\
 &\quad\times {}_4F_3\left(
 \begin{matrix}T_1+i\epsilon\lambda,\ldots,T_4+i\epsilon\lambda\\
 S/2+i\epsilon\lambda+ix,\ S/2+i\epsilon\lambda-ix,\
 1+2i\epsilon\lambda\end{matrix};1\right).
 \label{F}
\end{align}
Apparent coincident poles are evaluated by their analytic limit.
This is Proposition 4.4 in the symmetric dual-parameter notation.
With $D(T)_j=S/2-T_j$, his duality and parameter permutation symmetry give
\begin{equation}
 F(T;\lambda,x)=F(DT;x,\lambda),\qquad D^2T=T.
 \label{duality}
\end{equation}
The reversal of the four entries in the usual dual-parameter formula is
irrelevant because \eqref{F} is symmetric in those entries. The function
is even in each spectral variable. No new Wilson-function identity
or unitarity theorem is being claimed.

\newpage
\section{An exact simultaneous-reflection formula for Omega}
For $H\in\{h,1-h\}$ and $Q\in\{h',1-h'\}$ set
\[
 Q_0=H+h_1+h_3-1,\qquad z=Q_0-Q,
\]
\[
 A=(H+h_1+h_2-1,\ H+h_3+h_4-1,\ H+h_3-h_4,\ H+h_1-h_2),
\]
\[
 T_j=A_j-z/2,\qquad \lambda=\frac{z}{2i},\qquad
 x=\frac{Q_0-H-z/2}{i}.
\]
Direct matching of the two LPRS series \cite{lprs} with \eqref{F} yields
\begin{equation}
 \Om(H,Q)=K(H,Q)F(T;\lambda,x),
 \quad
 K=\Gamma(2H)\Gamma(2Q)\Gamma(A_1)\Gamma(A_2)
                    \Gamma(A_3-z)\Gamma(A_4-z).
 \label{dictionary}
\end{equation}
Indeed $S/2+i\lambda+ix=2Q_0-z$,
$S/2+i\lambda-ix=2H$, and $2i\lambda=z$.
The positive-$\lambda$ term gives the LPRS $\Gamma(-z)$ term,
and the negative term gives its $\Gamma(z)$ partner with the same $K$.
This verifies both coefficients, not merely the hypergeometric parameters.

Use a subscript $s$ for the data at $(1-H,1-Q)$, keeping the external
weights fixed. On the real scalar slice, exact algebra gives
\[
 \lambda_s=x,\qquad x_s=\lambda,\qquad DT=\overline{T_s}.
\]
Both $x$ and $\lambda$ are real. Evenness, conjugation and \eqref{duality}
therefore imply
\begin{equation}
 \boxed{\Om(H,Q)=R(H,Q)\,
             \overline{\Om(1-H,1-Q)}} ,
 \qquad R(H,Q)=\frac{K(H,Q)}{\overline{K(1-H,1-Q)}} .
 \label{reflection}
\end{equation}
Canceling the two common Gamma factors reduces this ratio to
\begin{align}
 R(H,Q)
 &=\frac{\Gamma(H+h_1+h_2-1)\Gamma(H+h_3+h_4-1)}
 {\Gamma(H+1-h_1-h_2)\Gamma(H+1-h_3-h_4)}
 \nonumber\\
 &\quad\times
 \frac{\Gamma(Q+1-h_1-h_4)\Gamma(Q+1-h_2-h_3)}
 {\Gamma(Q+h_1+h_4-1)\Gamma(Q+h_2+h_3-1)} .
 \label{R}
\end{align}
Every argument in \eqref{R} has real part $1/2$, so this ratio is finite
and nonzero on the slice. The derivation first avoids the apparent
Gamma-pair collisions and extends to them by continuity.

This uses analytic continuation of known special-function identities.
The complex parameter quadruples need not satisfy the positive-measure
domain of the Wilson transform; no Wilson spectral measure is substituted
for the physical $dS_3$ measure.

\newpage
\section{Complete-kernel stabilization and an absolute tail}
Retain the two physical coefficients $C_+$ and $C_-$, including the
phase of the earlier CG-period normalization, and put
\[
 X=\Om(h,h'),\qquad Y=\Om(1-h,h'),\qquad P=XY.
\]
Applying \eqref{reflection} to both factors of the negative-shadow term gives
\begin{equation}
 U_{00}=C_+P+\widehat C_-\,\overline P,\qquad
 \widehat C_-=C_-R(h,1-h')R(1-h,1-h').
 \label{stable}
\end{equation}
We evaluate the combined coefficient using logarithmic Gamma ratios.
There is no projection onto real values and no empirical symmetrization.

The modulus identity $|\Gamma(1/2+it)|^2=\pi/\cosh(\pi t)$ implies
\begin{equation}
 |R(h,1-h')R(1-h,1-h')|=
 \frac{\cosh\frac{\pi(d-a-p)}2\,\cosh\frac{\pi(d-q-c)}2}
 {\cosh\frac{\pi(d+a+p)}2\,\cosh\frac{\pi(d+q+c)}2}.
 \label{Rmod}
\end{equation}
In the product the two $H$-dependent Gamma ratios have combined
modulus one. The remaining two repeated $Q$ ratios give \eqref{Rmod}.
The preceding note's exact coefficient formula has precisely the
inverse ratio. Hence, for all positive labels,
\begin{equation}
 \boxed{|\widehat C_-|=|C_+|.}
 \label{balance}
\end{equation}
The old $e^{2\pi M}$ amplification is thus canceled analytically,
before a small numerical factor has to be computed.

Let $\widehat X,\widehat Y$ denote finite sums with exact-arithmetic
tail bounds $e_X,e_Y$ as in our preceding polynomial-certificate note.
Then
\begin{equation}
 |U_{00}-\widehat U_{00}|
 \le (|C_+|+|\widehat C_-|)
 (|\widehat X|e_Y+|\widehat Y|e_X+e_Xe_Y).
 \label{error}
\end{equation}
This is a bound for the complete scalar matrix element. It still excludes
rounding of finite sums, Gamma functions and coefficients.
It is an absolute bound; dividing by a nearly canceled final value is
not justified.

For the fixed-transfer collision path
\[
 c=M,\quad b=M+u,\quad a=M+u+v,\quad d=M+v,
\]
the coefficient modulus stays finite and nonzero:
\[
 |C_+|\longrightarrow
 \frac{e^{\pi(u-v+p+q)/4}}{4\pi^2}
 \sqrt{\cosh\frac{\pi(u+p)}2\cosh\frac{\pi(v-q)}2}.
\]
This removes one obstruction to the heavy-mass limit.
It does not supply uniform control of $X,Y$ or the physical mass
integrals. Low-mass direct sums can still converge slowly, and final
phase cancellation remains a separate source of relative sensitivity.

\newpage
\section{Checks, retained failures, and the next limit}
The exact symbolic dictionary is verified for six independent real labels.
Twenty-four checks use the previously saved four original Omega factors
at $M=2,4,8,16$ and offsets $r=0,0.1,0.4$.
Four fresh generic-label comparisons use the original LPRS formula,
and four compare \eqref{R} to the unsimplified $K/\overline K_s$.
Two independent Euler product integrals give a relative reflection
residual $4.75\,10^{-19}$ at 35 decimal digits.

The initial 33 checks retain two failures: the old $M=16$ negative-shadow
values at $r=0.1,0.4$ differ from their reflected counterparts by
$2.04\,10^{-60}$ and $7.42\,10^{-60}$, exceeding the preset $10^{-60}$
threshold. We did not relax that threshold or overwrite those values.
Fresh original-formula calculations at 100 and 130 digits reduce the
discrepancies to about $10^{-80}$ and $10^{-110}$, respectively.
The old exact-collision prototype bias is also retained and explicitly
separated from these checks.

For $p=q=0.6$, $u=0.3$, $v=-0.2$, we perform 18 complete-kernel runs:
six masses through $M=128$, cutoffs 64 and 128, and 70/110-digit refinement.
For $N=128$ the comparison with saved independent original-formula
references is:
\begin{center}
\begin{tabular}{rrrr}
\toprule
$M$ & $U_{00}$ (displayed real part) & reference discrepancy & tail bound\\
\midrule
8  &1.36981861885&$1.44\,10^{-13}$&$9.01\,10^{-12}$\\
16 &1.58731944205&$2.80\,10^{-23}$&$1.11\,10^{-21}$\\
32 &1.80680784696&$1.66\,10^{-42}$&$1.13\,10^{-40}$\\
\bottomrule
\end{tabular}
\end{center}
Full complex values are saved. The references themselves are not interval
enclosures. At $M=64,128$, no independent original-formula reference is
asserted. In particular, the $M=128$ analytic tail is $2.99\,10^{-127}$,
whereas changing arithmetic from 70 to 110 digits changes the value by
$5.35\,10^{-68}$. The tail is not the total numerical error.

Independent rational-polynomial and Decimal replay passes 365 checks,
including the original two failures, their four refinements, complex
residuals, coefficient moduli and propagated tails. Source/input hashes,
versions, and per-run checkpoints are saved.
The earlier collision stage's separate 100-digit parameter-derivative
job is still running and is not counted as passed here.

The next task is a uniform asymptotic analysis of the two remaining
positive-shadow factors, retaining their phases and fixed-versus-moving
preparation boundaries, followed by integration against the actual
creation amplitudes. Pointwise stabilization alone proves neither a
delta kernel nor the Bose free-observer limit. Full Fock domains,
finite-$G_N$ dynamics and the microscopic holographic dictionary remain open.

\begin{thebibliography}{9}
\bibitem{groenevelt} W. Groenevelt, \emph{The Wilson function transform},
\href{https://arxiv.org/html/math/0306424v1}{arXiv:math/0306424v1},
Sec.\ 3, Proposition 4.4 and Remark 4.5(iii).
\bibitem{lprs} J. Liu, E. Perlmutter, V. Rosenhaus and D. Simmons-Duffin,
\emph{$d$-dimensional SYK, AdS Loops, and $6j$ Symbols},
\href{https://arxiv.org/html/1808.00612v3}{arXiv:1808.00612v3},
Eqs.\ (3.35)--(3.38).
\end{thebibliography}
\end{document}
